% Complete copied proof fragment; source and changes are in BUILD_MANIFEST.json.
\section{Finite dilation of the original arithmetic packet over the tau
base}\label{proof04:finite-dilation-of-the-original-arithmetic-packet-over-the-tau-base}

\subsection{0. Scope and exact source
correspondence}\label{proof04:scope-and-exact-source-correspondence}

This calculation continues the original two-leg theta complex. The
retained infinitesimal identity is (A11) of
\texttt{arithmetic\_\allowbreak{}input.tex}; the absolute-base pushforward, its
two-leg action, and its tensor differential are (19), (21), (34), and
(46) of \emph{Arithmetic cohomology over the tau-base}. The full global
retraction and its cocycle were already constructed in transcript turn
\textbf{A1694}, node \texttt{1fa0ebbc-\allowbreak{}f569-\allowbreak{}40c1-\allowbreak{}b7ac-\allowbreak{}b523da19a5f0},
chain position 3576, and in \emph{Tau\_Global\_Comparison}, equations
(16), (19), (29)--(35), and (41). Those results are retained as prior
work, not reported as newly discovered here.

The additional calculation is the explicit finite-packet dilation
cocycle for every positive real dilation, its full Taylor-unit and
nilpotent coefficients, convergent source-space integral,
incomplete-gamma evaluation, exact finite/global change of
representatives, two-leg homotopies, and tensor and boundary formulas.
These formulas perform the group-level continuation of the existing
infinitesimal extension. They do not assign positivity to a
representative pairing.

The source bytes used are pinned at the end. Every integral on the
positive half-line keeps the displayed measure. In particular the
Hilbert pairing below uses \textbf{\(dx\)}, and its dilation mass is
\(a\), not \(1\). The notation \(t=\log a\) is a real time parameter;
the polynomial coordinate in \(\mathbb C[X]/(h)\) is written \(X\) to
avoid confusing it with that time. Thus \(A=M_X\) here is exactly the
source operator denoted \(M_t\).

\subsection{1. The original spaces, maps, and
topology}\label{proof04:the-original-spaces-maps-and-topology}

Work over \(\mathbb C\). Let

\[
V=\left\{\phi\in\mathcal S(\mathbb R):
\phi(-x)=\phi(x),\quad \phi(0)=0,\quad
\int_{\mathbb R}\phi(x)\,dx=0\right\},
\]

\[
\mathscr B=\left\{F\in C^\infty(\mathbb R_{>0}):
p_{b,k}(F)=\sup_{x>0}x^b|D^kF(x)|<\infty
\quad(b\in\mathbb Z,\ k\geq0)\right\},\qquad D=-x\partial_x.
\]

Use the Schwartz seminorms

\[q_{N,k}(\phi)=\sup_{x\in\mathbb R}(1+|x|)^N|\partial_x^k\phi(x)|,
\qquad N,k\geq0.\]

The space \(V\) is a closed subspace of the complete Schwartz space:
parity, value at zero, and integration are continuous. For example
\(|\int\phi|\leq q_{2,0}(\phi)\int_{\mathbb R}(1+|x|)^{-2}dx\). The
space \(\mathscr B\) is complete because a Cauchy sequence in all
\(p_{b,k}\) converges with all derivatives on compact subsets after
\(x=e^y\); the weighted uniform bounds pass to its limit. These
observations justify limits of Riemann sums in the actual spaces below.

Retain

\[
\widehat\phi(\xi)=\int_{\mathbb R}\phi(x)e^{-2\pi ix\xi}\,dx,
\quad \Theta\phi(x)=2\sum_{n\geq1}\phi(nx),
\quad JF(x)=x^{-1}F(1/x),
\]

\[\mathcal MF(s)=\int_0^\infty F(x)x^s\,\frac{dx}{x}.\]

Schwartz estimates, followed by the original Poisson identity, give

\[\Theta:V\longrightarrow\mathscr B,\qquad
\Theta\widehat\phi=J\Theta\phi,\qquad D\Theta=\Theta D.\]

Here is the continuity argument. At \(x\geq1\), each \(D^k\phi(nx)\) is
bounded by a finite sum of Schwartz seminorms times \((nx)^{-M}\), where
\(M\) can be any sufficiently large positive integer. Summing \(n^{-M}\)
gives every required bound at infinity. The Fourier transform preserves
Schwartz seminorms continuously, preserves evenness, and interchanges
the two vanishing moments. Poisson summation therefore gives
\(\Theta\phi(x)=x^{-1}\Theta\widehat\phi(1/x)\); applying \(D^k\) and
the preceding bounds gives every required estimate at zero. Each
resulting estimate involves finitely many source seminorms. Termwise
differentiation is justified uniformly on every \(x\geq\delta>0\) by the
same summable majorants.

The theta map is injective. For fixed \(x>0\), absolute convergence
permits

\[
\frac12\sum_{m\geq1}\mu(m)\Theta\phi(mx)
=\sum_{r\geq1}\left(\sum_{m\mid r}\mu(m)\right)\phi(rx)=\phi(x).
\]

Indeed the rearranged absolute sum is bounded by
\(\sum_{r\geq1}d(r)|\phi(rx)|<\infty\), since \(d(r)\leq r\) and
\(\phi\) is Schwartz. Thus no theta relation is lost when its source is
recovered.

For \(a>0\) define \(U_aF(x)=F(x/a)\), and use the same formula on
\(V\). Direct substitutions give

\[
p_{b,k}(U_aF)=a^b p_{b,k}(F),\qquad
q_{N,k}(U_a\phi)\leq a^{-k}\max(1,a)^Nq_{N,k}(\phi),
\tag{1}
\]

\[
\mathcal MU_aF(s)=a^s\mathcal MF(s),\quad
\widehat{U_a\phi}=aU_{1/a}\widehat\phi,\quad
\Theta U_a=U_a\Theta.\tag{2}
\]

The two moments of \(U_a\phi\) are \(\phi(0)\) and \(a\int\phi\), so
\(V\) is invariant. Differentiating the defining formula gives

\[\frac{d}{dt}U_{e^t}\phi=DU_{e^t}\phi=U_{e^t}D\phi.\]

This is a derivative in every Schwartz seminorm. To verify it, apply the
ordinary Taylor formula with integral remainder to \(U_{e^{t+h}}\phi\);
its second derivative is \(U_{e^{t+h}}D^2\phi\), uniformly bounded on
each compact time interval by (1). The same argument, using the exact
first identity in (1), proves smoothness in \(\mathscr B\). Repeating
with \(D^j\phi\) proves derivatives of all orders. In particular this is
a smooth group action on the retained Frechet spaces.

The four-point source poset is

\[+<\eta<\sigma,\qquad -<\eta<\sigma.\]

For the sheaf with stalks \(V,V,\mathscr B,0\) and restriction maps
\(\Theta,\Theta\mathcal F\), the source projective resolution computes

\[
C=\mathsf A_\tau(\mathcal T)
=\left[V\oplus V\xrightarrow{d}\mathscr B\right],
\quad d(v,w)=\Theta(v-\widehat w),\quad |V\oplus V|=0,\ |\mathscr B|=1.
\tag{3}
\]

Its actual chain action is

\[
\mathcal U_a^0(v,w)=(U_av,aU_{1/a}w),\qquad
\mathcal U_a^1F=U_aF.\tag{4}
\]

Equation (2) proves \(d\mathcal U_a^0=\mathcal U_a^1d\). On even
Schwartz functions \(\mathcal F^2=1\), so the exact degree-zero cycle
inclusion is

\[i_0(w)=(\widehat w,w),\qquad
\mathcal U_ai_0(w)=i_0(aU_{1/a}w).\tag{5}\]

It follows from theta injectivity that \(H^0(C)=i_0(V)\) and
\(H^1(C)=Q=\mathscr B/\Theta V\). These are the original global
cohomology groups over the absolute tau base.

\subsection{2. The finite arithmetic packet and its unchanged Taylor
unit}\label{proof04:the-finite-arithmetic-packet-and-its-unchanged-taylor-unit}

Retain the original Gaussian source

\[\phi_*(x)=(4\pi^2x^4-6\pi x^2)e^{-\pi x^2},\qquad f_0=\Theta\phi_*.\]

It has value zero at zero, and direct Gaussian integration gives zero
integral. More precisely, with
\(M_+\phi(s)=\int_0^\infty\phi(x)x^s\,dx/x\),

\[
M_+\phi_*(s)=\pi^{-s/2}\bigl(2\Gamma(s/2+2)-3\Gamma(s/2+1)\bigr)
=\tfrac12s(s-1)\pi^{-s/2}\Gamma(s/2).\tag{6}
\]

Consequently the entire function \(g=\mathcal Mf_0\) agrees, first by
absolutely convergent integration in a right half-plane and then by
continuation, with the original function

\[g(s)=2\xi(s)=s(s-1)\pi^{-s/2}\Gamma(s/2)\zeta(s).\]

Entireness of the Mellin transform of \(f_0\in\mathscr B\) follows
directly from the bounds in \(\mathscr B\). For
\(0<\operatorname{Re}s<1\), set

\[H_\phi(s)=\frac{2\pi^{s/2}}{s(s-1)\Gamma(s/2)}M_+\phi(s).\]

Then

\[\mathcal M\Theta\phi=gH_\phi,\qquad H_{\phi_*}=1,\qquad
H_{P(D)\phi_*}(s)=P(s).\tag{7}\]

For completeness, \(M_+\phi\) is holomorphic on
\(\operatorname{Re}s>-2\), since smooth evenness and \(\phi(0)=0\) give
\(\phi(x)=O(x^2)\). In \(\operatorname{Re}s>1\), interchanging its theta
series and Mellin integral gives
\(\mathcal M\Theta\phi=2\zeta(s)M_+\phi(s)\). Continue to the indicated
strip to obtain the first equality in (7). Integrating \(D\phi\) by
parts proves \(H_{D\phi}=sH_\phi\); the endpoint term
\([x^s\phi(x)]_0^\infty\) vanishes on that strip. Equation (6) then
proves the other two equalities.

Fix the original nonempty finite packet \(Z\) of actual nontrivial
zeros, each with its complete order \(m_\rho\), and put

\[
h(X)=\prod_{\rho\in Z}(X-\rho)^{m_\rho},\quad d_h=\deg h,
\quad E=\mathbb C[X]/(h),\quad A=M_X,
\]

\[v(s)=g(s)/h(s),\qquad \varepsilon=j_h(h/g)\in E^\times.\]

The quotient \(h/g\) here is taken as a germ at each selected zero;
these germs are holomorphic and nonzero because the complete
multiplicity has been retained. The local unit is \textbf{not} replaced
by \(1\). For \(u\in E\), let \(R_Z(u)\) be the unique polynomial of
degree less than \(d_h\) representing \(\varepsilon u\), and put

\[\ell(u)=[X^{d_h-1}]R_Z(u).\tag{8}\]

We record the inverse used to define the source section, rather than
taking an unproved inverse as an extra hypothesis. On \(F\in\mathscr B\)
satisfying \(\mathcal MF(\rho)=0\), define

\[S_\rho F(x)=x^{-\rho}\int_x^\infty F(y)y^{\rho-1}\,dy.\]

Differentiation gives \((D-\rho)S_\rho F=F\). The moment condition makes
the upper integral the negative lower integral. If
\(\sigma=\operatorname{Re}\rho\), the upper estimate, for \(N>\sigma\),
is

\[|S_\rho F(x)|\leq \frac{p_{N,0}(F)}{N-\sigma}x^{-N}.\]

The lower estimate, for \(N+\sigma>0\), is

\[|S_\rho F(x)|\leq \frac{p_{-N,0}(F)}{N+\sigma}x^{N}.\]

Using the upper estimate on \(x\geq1\) and the lower estimate on
\(x\leq1\) gives every weighted bound. The differential identity
expresses \(D^kS_\rho F\) as \(\rho^kS_\rho F\) plus a finite linear
combination of \(D^jF\), proving all derivative bounds and continuity.
The homogeneous solution \(cx^{-\rho}\) cannot belong to \(\mathscr B\)
unless \(c=0\), so the inverse is unique. Integration by parts gives

\[\mathcal MS_\rho F(s)=\frac{\mathcal MF(s)}{s-\rho},\]

with the removable value filled in. Iterating this construction,
removing one vanishing factor at a time, defines the continuous inverse
\(S_h\) of \(h(D)\) on \(\ker(j_h\mathcal M)\). The operation preserves
the remaining zeros and lowers the order at the chosen zero by exactly
one at each step. Conversely any \(h(D)F\) has those vanishing jets.
This proves

\[h(D)\mathscr B=\ker(j_h\mathcal M),\qquad h(D)\text{ is injective}.\tag{9}\]

The section in the question is therefore the actual well-defined map

\[R=R_{\rm ref}:E\longrightarrow\mathscr B,\qquad
Ru=S_h\Theta\bigl(R_Z(u)(D)\phi_*\bigr).\tag{10}\]

Indeed its input has Mellin transform \(gR_Z(u)\), which vanishes to all
the required orders. Equations (7) and (9) give

\[\mathcal MRu=vR_Z(u),\qquad
J_ZR=1_E,\qquad J_Z=j_h\mathcal M,\qquad J_Z\Theta=0.\tag{11}\]

The Mellin-jet map is continuous: split its integral at \(1\), dominate
the logarithmic factors at zero by a sufficiently negative weight and
those at infinity by a sufficiently positive weight. This yields a
finite sum of seminorm bounds for each derivative at each \(\rho\). In
(11), \(j_h(v)\varepsilon=1\) proves the middle identity. Hence \(R\)
and \(\sigma=qR:E\to Q\) are injective, since \(J_Z\) descends to \(Q\).

The infinitesimal boundary is precisely the retained one:

\[DR-RA=\Theta\phi_*\ell=f_0\ell.\tag{12}\]

To prove it without losing its coefficient, the polynomial
\(XR_Z(u)-R_Z(Au)\) vanishes modulo \(h\) and has degree at most
\(d_h\). Since \(h\) is monic, its quotient by \(h\) is exactly
\(\ell(u)\). Apply \(h(D)\) to (12), use (10), \(D\Theta=\Theta D\), and
that polynomial identity. Injectivity of \(h(D)\) from (9) proves (12).
For a single simple zero, \(R_Z(u)=u/g'(\rho)\) and
\(\ell(u)=u/(2\xi'(\rho))\), with the original factor of two.

\subsection{3. The finite dilation cocycle as an actual convergent
source
integral}\label{proof04:the-finite-dilation-cocycle-as-an-actual-convergent-source-integral}

For every real \(t\), define \(T_t=e^{tA}\) and

\[
c_t(u)=\int_0^t U_{e^{t-r}}\phi_*\,\ell(T_ru)\,dr\in V.
\tag{13}
\]

If \(t<0\), \(\int_0^t=-\int_t^0\); no orientation is changed. Write
\(c_a=c_{\log a}\) and \(a^A=T_{\log a}\) when using the multiplicative
parameter.

\textbf{Convergence and continuity.} Fix any norm on the
finite-dimensional coefficient space, solely for estimates, and its
induced operator norms. On the oriented interval between \(0\) and
\(t\), both \(|r|\) and \(|t-r|\) are at most \(|t|\). Equation (1) and
\(\|T_r\|\leq e^{|r|\|A\|}\) yield

\[
q_{N,k}(c_tu)
\leq |t|e^{(N+k+\|A\|)|t|}\,
q_{N,k}(\phi_*)\,\|\ell\|\,\|u\|.\tag{14}
\]

The integrand is continuous in every seminorm. Its Riemann sums are
Cauchy in every seminorm by uniform continuity on the compact interval.
Completeness of \(V\) proves their common limit exists in \(V\), and the
inequality follows first for sums and then for their limit. Thus the
integral is an actual source vector with the two original zero moments
and even Schwartz regularity. It is not merely an integral in a weak
dual or in \(L^2\).

Smoothness in \(t\) follows either by the differentiable
compact-interval rule justified by the same estimates for all time
derivatives, or by writing (13) as
\(t\int_0^1 U_{e^{t(1-r)}}\phi_*\ell(T_{tr}u)dr\). This latter formula
works without a sign convention change for negative \(t\). Every
derivative has a finite sum of terms involving \(D^j\phi_*\), powers of
\(A\), and bounded polynomial factors in \(r\). Inequality (14) with
those finitely many seminorms proves local uniform convergence of the
differentiated Riemann sums. Therefore \(t\mapsto c_t\) is smooth as a
map to continuous linear maps \(E\to V\), with their bounded-set
topology. Finite dimensionality of \(E\) makes the preceding uniform
estimates on its unit ball sufficient for this statement.

\textbf{The exact dilation identity.} For all \(t\in\mathbb R\),

\[\boxed{U_{e^t}R-RT_t=\Theta c_t.}\tag{15}\]

For fixed \(u\), differentiate in the Frechet space \(\mathscr B\):

\[
\frac{d}{dr}\bigl(U_{e^{t-r}}RT_ru\bigr)
=-U_{e^{t-r}}(DR-RA)T_ru
=-\Theta U_{e^{t-r}}\phi_*\ell(T_ru).
\]

The fundamental theorem of calculus follows in every defining seminorm
from the same Riemann-sum argument; integrating with the indicated
orientation proves (15), including \(t<0\). Continuity of \(\Theta\)
permits moving it through the integral.

\textbf{Cocycle and derivative.} Directly splitting the defect in (15)
gives

\[
\boxed{c_{t+s}=U_{e^t}c_s+c_tT_s,\qquad
c_0=0,\qquad
c_{-t}=-U_{e^{-t}}c_tT_{-t}.}\tag{16}
\]

Indeed applying \(\Theta\) to the asserted first equality yields the
identity obtained by inserting \(U_{e^t}RT_s\) in
\(U_{e^{t+s}}R-RT_{t+s}\). Injectivity of \(\Theta\) proves equality in
the actual source \(V\). The other two statements follow by setting
\(t=s=0\) and \(s=-t\). Differentiating (13), or the cocycle in either
variable, gives the two equal expressions

\[
\boxed{\partial_tc_t=Dc_t+\phi_*\ell T_t
=U_{e^t}\phi_*\ell+c_tA,\qquad
\partial_tc_t|_{t=0}=\phi_*\ell.}\tag{17}
\]

All derivatives are in \(V\). In multiplicative notation (16) is

\[c_{ab}=U_a\circ c_b+c_a\circ b^A,\]

where the final product means the map \(c_a\circ b^A\); to remove any
typographical ambiguity it can also be written
\(c_{ab}(u)=U_a(c_b(u))+c_a(b^Au)\).

The space \(\mathscr B_Z=\Theta V+RE\) is a direct topological sum with
coordinates \((\phi,u)\mapsto\Theta\phi+Ru\). Its inverse on that
subspace is \((\Theta^{-1}(F-RJ_ZF),J_ZF)\); continuity of the inverse
theta map is also proved in Section 6 below. Equations (15) and (16)
integrate the source block operator to the exact group representation

\[
U_a\big|_{\mathscr B_Z}
=\begin{pmatrix}U_a|_V&c_a\\0&a^A\end{pmatrix}.\tag{18}
\]

This is the explicit group-level form of the already known extension
(A12). In particular (15) proves the true quotient intertwining
\(\overline U_a\sigma=\sigma a^A\).

\subsection{4. All nilpotents, the Mellin multiplier, and an
incomplete-gamma
evaluation}\label{proof04:all-nilpotents-the-mellin-multiplier-and-an-incomplete-gamma-evaluation}

The Chinese remainder decomposition is

\[E=\bigoplus_{\rho\in Z}\mathbb C[z_\rho]/(z_\rho^{m_\rho}),\qquad
A|_\rho=\rho I+N_\rho,\]

where \(N_\rho z_\rho^j=z_\rho^{j+1}\) and \(N_\rho^{m_\rho}=0\). Denote
the summand projector by \(P_\rho\), extending \(N_\rho\) by zero on
other summands. These are the actual full local algebras. They can be
constructed explicitly: if \(H_\rho(X)=h(X)/(X-\rho)^{m_\rho}\), then
\(P_\rho\) is multiplication by the remainder modulo \(h\) of

\[H_\rho(X)\sum_{j=0}^{m_\rho-1}
\frac{(H_\rho^{-1})^{(j)}(\rho)}{j!}(X-\rho)^j.\]

At every other centre this has all required zero jets; at \(\rho\) it
has jet \(1\). This verifies all its projector identities by the Chinese
remainder map.

The action and the entire boundary integral therefore contain the full
expressions

\[
T_t=\sum_{\rho\in Z}e^{t\rho}
\sum_{j=0}^{m_\rho-1}\frac{t^j}{j!}N_\rho^jP_\rho,\tag{19}
\]

\[
c_tu=\sum_{\rho\in Z}\sum_{j=0}^{m_\rho-1}
\ell(N_\rho^jP_\rho u)
\int_0^t U_{e^{t-r}}\phi_*\,e^{r\rho}\frac{r^j}{j!}\,dr.
\tag{20}
\]

The unit entering every coefficient is still \(\varepsilon=j_h(h/g)\)
through (8). If the packet has one zero of order \(m\), and
\(\varepsilon=\sum_{i<m}\varepsilon_i z^i\), \(u=\sum_{i<m}u_i z^i\),
then \(\ell(u)=\sum_{i=0}^{m-1}\varepsilon_i u_{m-1-i}\). This follows
because the top coefficient of a polynomial of degree less than \(m\) is
the same in the \(X\) and \(z=X-\rho\) bases. Thus even the functional
\(\ell\) retains the lower coefficients of the Taylor unit, rather than
only its nonzero constant value.

There is a closed expression for the integral kernels in (20). For
\(x>0\) let

\[
K_{\rho,t}(x)=\int_0^t e^{r\rho}\phi_*(xe^{r-t})\,dr.
\]

Define the oriented incomplete-gamma interval by

\[\Gamma_{[L,H]}(w)=\int_L^H y^{w-1}e^{-y}\,dy\qquad(L,H>0),\]

where \(y^{w-1}=\exp((w-1)\log y)\) uses the real logarithm. Then
substitution first of \(y=xe^{r-t}\) and then of \(z=\pi y^2\) gives

\[
\begin{aligned}
K_{\rho,t}(x)
={}& e^{t\rho}x^{-\rho}\pi^{-\rho/2}
\left[
2\Gamma_{[\pi x^2e^{-2t},\,\pi x^2]}(\rho/2+2)
-3\Gamma_{[\pi x^2e^{-2t},\,\pi x^2]}(\rho/2+1)
\right].
\end{aligned}\tag{21}
\]

Both substitutions preserve orientation when \(t<0\). The coefficients
\(2\) and \(-3\) come respectively from \(4\pi^2x^4\) and \(-6\pi x^2\),
including \(dy/y=dz/(2z)\). No Gaussian constant is removed.

Differentiation under the compact oriented \(r\)-integral now gives the
finite formula

\[\boxed{c_tu=\sum_{\rho\in Z}\sum_{j=0}^{m_\rho-1}
\frac{\ell(N_\rho^jP_\rho u)}{j!}\,
\partial_\rho^jK_{\rho,t}.}\tag{22}\]

The derivative here acts on the complex parameter in the entire kernel
\(K_{\rho,t}\) before evaluation at the indicated centre, not on the
arithmetic zero or on the coefficient \(\ell(N_\rho^jP_\rho u)\). The
formula on \(x>0\) is extended evenly to \(\mathbb R\), with its smooth
value at \(0\) supplied by the integral defining \(K\). That integral
and Section 3 prove smoothness and both zero moments, so no regularity
is inferred from cancelling the separate singular-looking factors in
(21).

The full Mellin multiplier of the source boundary is also explicit:

\[
\boxed{H_{c_tu}(s)=\int_0^t e^{(t-r)s}\ell(T_ru)\,dr
=\frac{e^{ts}R_Z(u)(s)-R_Z(T_tu)(s)}{h(s)}.}\tag{23}
\]

For the first equality in the original strip, use (7),
\(H_{U_a\phi_*}=a^s\), and the continuous source integral. For the
second, differentiate \(e^{(t-r)s}R_Z(T_ru)(s)\) in \(r\). The
polynomial identity proving (12) says its derivative is
\(-e^{(t-r)s}h(s)\ell(T_ru)\). Integration proves (23) wherever
\(h(s)\ne0\). The numerator has every required vanishing jet: in the
local algebra its jet is \(e^{tA}\varepsilon u-\varepsilon T_tu=0\).
Hence the quotient in (23) has removable singularities at every chosen
zero and is entire. This entire expression extends the original source
multiplier; it is not a replacement of \(h\) or \(g\).

For every integer \(n\geq0\) its actual Taylor coefficient at any
complex centre \(s_0\) is

\[
\frac{H_{c_tu}^{(n)}(s_0)}{n!}
=\int_0^t e^{(t-r)s_0}\frac{(t-r)^n}{n!}\ell(T_ru)\,dr.
\tag{24}
\]

Inserting (19) gives every mixed coefficient in \((t-r)^n r^j/(n!j!)\)
and keeps all nilpotent terms. This is also a direct convergence proof
for every jet, since the parameter interval is compact.

\subsection{5. Actual two-leg homotopies and the support-face
comparison}\label{proof04:actual-two-leg-homotopies-and-the-support-face-comparison}

Regard \(E[-1]\) as the complex with \(E\) in degree \(1\) and zero
elsewhere. The finite section defines a degree-zero cochain map
\(r:E[-1]\to C\), with \(r^1=R\). The maps \(\mathcal U_a r\) and
\(r a^A\) have homotopies

\[
H_a^{+}:E[-1]\longrightarrow C[-1],\quad
(H_a^{+})^1u=(c_au,0),
\]

\[
H_a^{-}:E[-1]\longrightarrow C[-1],\quad
(H_a^{-})^1u=(0,-\widehat{c_au}).\tag{25}
\]

Here the notation records maps of cohomological degree \(-1\); the
source only has degree \(1\). In ordinary component notation the exact
identity is

\[\boxed{\mathcal U_ar-ra^A=dH_a^\pm+H_a^\pm d_{E[-1]}=dH_a^\pm.}\tag{26}\]

Indeed \(d(c_au,0)=\Theta c_au\), whereas
\(d(0,-\widehat{c_au})=\Theta\mathcal F^2c_au=\Theta c_au\). This
verifies the minus sign using the original differential. Equation (2)
also gives the homotopy cocycle itself:

\[H_{ab}^\pm=\mathcal U_a^0H_b^\pm+H_a^\pm b^A.\tag{27}\]

No extra Fourier scalar is inserted:
\(\widehat{U_ac}=aU_{1/a}\widehat c\) is exactly the second-leg action
in (4).

The original support masks are \(\lambda\in\{\{+\},\{-\},\{+,-\}\}\).
The fibre \(C_\lambda^0\) consists of the legs present in \(\lambda\),
\(C_\lambda^1=\mathscr B\), with the same differential. The \(+\)
homotopy in (25) is defined on \(\{+\}\) and on the joint face. The
\(-\) homotopy is defined on \(\{-\}\) and on the joint face. On the
joint face their difference is the precisely retained cycle

\[H_a^+-H_a^-=(c_a,\widehat c_a)=i_0(\widehat c_a).\tag{28}\]

It lies in \(H^0(C)\) by (5) and theta injectivity. There is no term in
degree \(-1\) of \(C\), so a nonzero such cycle is not a boundary in
\(C\).

The two singleton homotopies are forced. On \(\{+\}\), injectivity of
theta forces the first one; on \(\{-\}\) it forces the second.
Consequently a single homotopy that is strictly natural for both
singleton-to-joint inclusions would require
\((c_a,0)=(0,-\widehat c_a)\), which is equivalent to \(c_a=0\). This is
the exact scope of the support-naturality obstruction; it does not
remove the homotopies in (25), their difference (28), or the resulting
quotient action.

For completeness, \(c_a\) as a linear map is nonzero for every
\(a\ne1\). If \(c_a=0\), (15) would make the nonzero finite-dimensional
subspace \(RE\) invariant under \(U_a\), with invertible action \(a^A\).
It therefore has a nonzero eigenvector \(F\in\mathscr B\) with
\(U_aF=\lambda F\), \(\lambda\ne0\). Replacing \(a\) by \(a^{-1}\) if
needed, take \(a>1\). Choose \(x_0>0\) with \(F(x_0)\ne0\). Iteration
gives \(F(a^nx_0)=\lambda^{-n}F(x_0)\) for \(n\geq0\). The bound
\(p_{N,0}(F)<\infty\) implies
\(a^{nN}|\lambda|^{-n}|F(x_0)|\leq x_0^{-N}p_{N,0}(F)\). Choose \(N\)
with \(a^N>|\lambda|\) and let \(n\to\infty\) to obtain a contradiction.
Thus no such finite invariant subspace exists. This extends the
previously proved failure of a strict equivariant section by giving its
exact nonzero finite-dilation defect; the failure itself was already
present in (A12).

At a fixed admitted mask the split lift of every displayed map is

\[\widetilde f(\lambda,u)=(\lambda,f(u)),\qquad \widetilde f(\tau)=\tau.\]

An amplitude \(f(u)=0\) has value \((\lambda,0)\), the supported zero,
including at \(a=1\). All homotopy identities are interpreted using
addition and the supported scalar \((-1)^\bullet\) in that fibre. No
absent leg is inserted by the face-specific formulas.

The known synchronization comparison can be retained explicitly when a
common joint target is desired. Send every active degree-one mask to
\((\{+,-\},F)\) and send the degree-zero pair to its joint-supported
pair with the same two amplitudes, putting a supported zero in a
formerly missing coordinate. Denote this operation by \(\mathfrak r\).
It is the source map \((a,b)\mapsto(a+eb,b+ea)\), because multiplication
by \(e\) gives a supported zero at a present coordinate. It preserves
the differential, fixes joint masks, and is idempotent. The pair
\((\operatorname{supp}x,\mathfrak r x)\) recovers the original element
by selecting the originally present coordinates, proving injectivity of
that comparison. After this explicit synchronization, (25)--(28) hold in
the common joint fibre. The original mask must stay attached to the
comparison; forgetting it is not part of the homotopy calculation.

One can also see the extra degree-zero group in an exact equivariant
roof. Let \(C_Z=[V\oplus V\to\mathscr B_Z]\) with the same differential,
and define

\[p_Z^0(v,w)=w,\qquad p_Z^1=J_Z.\]

This is a cochain map to \(V[0]\oplus E[-1]\). It is equivariant for the
action \(aU_{1/a}\) on \(V\) and \(a^A\) on \(E\), and induces
isomorphisms on both cohomology groups: the degree-zero assertion
follows from (5), while the degree-one assertion follows from the direct
sum \(\mathscr B_Z=\Theta V\oplus RE\) and (11). Thus the exact two-leg
version of the source roof is

\[V[0]\oplus E[-1]\ \xleftarrow{\ p_Z\ }\ C_Z\ \longrightarrow\ C.\]

Its ordinary cochain section is \(i_Z^0(w)=(\widehat w,w)\) and
\(i_Z^1=R\). If \(F=\Theta\phi+Ru\), the map \(h_Z^1(F)=(\phi,0)\) obeys
\(1-i_Zp_Z=dh_Z+h_Zd\): in degree zero it returns \((v-\widehat w,0)\)
and in degree one it returns \(\Theta\phi\). Its failure of equivariance
in degree one is exactly (25). This proves the replacement on the entire
original two-leg complex, retaining its \(V[0]\) summand.

\subsubsection{The transported minus face has an explicitly nonzero
extension
component}\label{proof04:the-transported-minus-face-has-an-explicitly-nonzero-extension-component}

The preceding roof also computes what happens to the original support
arrows when the action is retained. Work over \(\mathbb C[X]\), with
\(X=D\) on \(\mathscr B_Z\), \(X=A\) on \(E\), and \(X=1-D\) on the
degree-zero second-leg module, denoted \(V_-\). Let \(r(X)=R_Z(1)\) and
\(\psi=r(D)\phi_*\). Then

\[h(D)R(1)=\Theta\psi.\]

The class of \(\psi\) in \(V/h(D)V\) is nonzero: applying \(j_hH\) gives
\(j_hr=\varepsilon\ne0\), whereas \(j_hH\) vanishes on \(h(D)V\) by (7).
The Fourier map carries this quotient isomorphically onto \(V/h(1-D)V\),
since differentiating (2) gives \(\mathcal F D=(1-D)\mathcal F\).

For the minus-face extension the injection is \(i_-=-\Theta\mathcal F\).
Its extension representative, with the convention \(h(D)R(1)=i_-(w)\),
is therefore

\[w=-\widehat\psi\quad\text{in }V/h(1-D)V.\]

The component of the minus support arrow in the \(V_-[0]\) summand of
the joint roof has the opposite sign to this extension representative.
Here is the exact shifted-resolution calculation. The free cochain
resolution of \(E\) in degrees \(-1,0\) has differential \(h\). We use
\((C[n])^j=C^{j+n}\) and \(d_{C[n]}=(-1)^nd_C\); after shifting by
\(-1\), that resolution is in degrees \(0,1\) with differential \(-h\).
Its degree-one lift to the minus-face complex sends \(1\) to \(R(1)\),
extended \(\mathbb C[X]\)-linearly. Its degree-zero lift must send \(1\)
to \(-w\), because \(i_-(-w)=-h(D)R(1)\). Projection to \(V_-[0]\) is
identity in degree zero. Thus the induced derived map
\(E[-1]\to V_-[0]\) represents exactly

\[
[-w]=[\widehat\psi]
\quad\text{in}\quad
\operatorname{Ext}^1_{\mathbb C[X]}(E,V_-)
\cong V/h(1-D)V.\tag{28a}
\]

The isomorphism here follows by applying
\(\operatorname{Hom}_{\mathbb C[X]}(-,V_-)\) to that two-term free
resolution; its cokernel is the stated quotient. The class (28a) is
nonzero by the preceding source-jet calculation. The same arrow has
identity as its \(E[-1]\) component. By contrast, the plus support arrow
has zero \(V_-[0]\) component and identity \(E[-1]\) component, because
its degree-zero image has second coordinate zero. Consequently the two
support arrows in the equivariant derived category do not both become
inclusion of the \(E[-1]\) summand: the exact difference carries (28a),
with the original Taylor unit. The strict original cochain support
diagram remains available before passage through the roof.

\subsection{6. The exact bridge to the already constructed global
section}\label{proof04:the-exact-bridge-to-the-already-constructed-global-section}

We recall the global section in a form that makes its use here
mathematically explicit. Choose the same fixed even real cutoffs
\(\alpha,\beta\in C_c^\infty(\mathbb R)\), with
\(0\leq\alpha,\beta\leq1\), equal to \(1\) near zero. For \(x\ne0\)
define

\[\mathcal UF(x)=\tfrac12\sum_{n\geq1}\mu(n)F(n|x|).\]

If \(P_j(D)=(-1)^jD(D+1)\cdots(D+j-1)=x^j\partial_x^j\), and
\(C_{N,j}(F)\) is the sum of \(p_{N,k}(F)\) weighted by the absolute
coefficients of \(P_j\), then for \(x>0\), \(N>1\),

\[|\partial_x^j\mathcal UF(x)|\leq\tfrac12\zeta(N)C_{N,j}(F)x^{-N-j}.\tag{29}\]

Each summand is bounded by \(\tfrac12C_{N,j}(F)n^{-N}x^{-N-j}\);
summation proves the estimate and local uniform convergence of every
derivative. Thus the cutoff exterior functions \((1-\alpha)\mathcal UF\)
and \((1-\beta)\mathcal UJF\), extended by zero near zero, are
continuous Schwartz-valued maps. The theta inversion in Section 1 and
Poisson give \(\mathcal U\Theta\phi=\phi\) and
\(\mathcal UJ\Theta\phi=\widehat\phi\) away from zero.

On \(L^2(\mathbb R,dx)\) define

\[T=M_\alpha\mathcal F^{-1}M_\beta\mathcal F,
\qquad
YF=(1-\alpha)\mathcal UF+
\alpha\mathcal F^{-1}((1-\beta)\mathcal UJF).\]

The kernel of \(T\) is \(\alpha(x)\check\beta(x-y)\); its squared
integral is \(\|\alpha\|_2^2\|\beta\|_2^2\), so \(T\) is compact. Its
norm \(c\) is strictly less than \(1\). Otherwise compactness would give
a unit vector attaining norm \(1\). Equality in the two contraction
inequalities forces its Fourier transform to be supported where
\(\beta=1\), and then the vector itself to be supported where
\(\alpha=1\). A compactly supported Fourier transform in \(L^2\) is in
\(L^1\), and its inverse Fourier transform extends to an entire function
by its finite-interval integral. That entire function would vanish on a
real open interval because of the compact spatial support, hence vanish
identically, a contradiction.

Thus \((1-T)^{-1}=\sum_{n\geq0}T^n\) in operator norm, with bound
\((1-c)^{-1}\). The operator \(T:L^2\to\mathcal S\) is continuous:
Cauchy--Schwarz on the compact Fourier support bounds every derivative
of \(\mathcal F^{-1}(\beta\widehat f)\) by a constant times \(\|f\|_2\),
and multiplication by \(\alpha\) gives compact spatial support and every
Schwartz seminorm. It follows that \(\Phi=(1-T)^{-1}Y\) is
Schwartz-valued and continuous; the explicit identity \(\Phi=Y+T\Phi\)
gives

\[q_{N,k}(\Phi F)\leq q_{N,k}(YF)+\frac{C_{N,k}}{1-c}\|YF\|_2.\]

No uniform lower bound for \(1-c\) as cutoffs vary is used. On a theta
input, \(Y\Theta\phi=(1-T)\phi\), so \(\Phi\Theta\phi=\phi\). The range
of \(\Phi\) is the even Schwartz space
\(\mathcal S(\mathbb R)^{\rm even}\). To check this exact domain,
\(\mathcal UF\) is even because its defining series uses \(|x|\), and
the same is true of \(\mathcal UJF\). The cutoffs \(\alpha,\beta\) are
even, and Fourier transformation and inverse Fourier transformation
preserve evenness, so \(YF\) is even. Reflection
\(\phi(x)\mapsto\phi(-x)\) commutes with \(T\); hence it commutes with
every partial sum of the norm-convergent inverse \((1-T)^{-1}\) and with
its limit. Thus \(\Phi F\) is even as an \(L^2\) vector, and its
Schwartz representative is even pointwise by continuity. With the
original Gaussian moment corrections, define the map on this precise
domain by

\[
P_V\phi=\phi-\phi(0)e^{-\pi x^2}
-\left(\int_{\mathbb R}\phi-\phi(0)\right)2\pi x^2e^{-\pi x^2},
\qquad \Lambda=P_V\Phi.
\]

The moment columns of the two displayed Gaussians are \((1,1)\) and
\((0,1)\) by direct Gaussian integration. The two subtractions therefore
set the value at zero and the integral to zero, while preserving
evenness. On \(V\) both correction coefficients vanish. Consequently
\(P_V:\mathcal S(\mathbb R)^{\rm even}\to V\) is a continuous projection
and \(P_V|_V=1\). This is its stated domain; no projection of the entire
Schwartz space onto the even subspace is asserted by that formula. It
follows that

\[\Lambda\Theta=1_V,\quad K=1-\Theta\Lambda,\quad
s:Q\to\mathscr B,\quad s[F]=KF.\tag{30}\]

This proves all needed existence and continuity claims for the
particular global construction used in A1694. Since \(\ker K=\Theta V\),
that image is closed. The expression for \(s\) is well-defined,
\(qs=1\), \(\Lambda s=0\), and it is continuous by the definition of the
quotient topology. In particular \(\Theta^{-1}\) on its image is
\(\Lambda\), verifying the topological inverse asserted after (17).

The already existing global cocycle is \(k_a=\Lambda U_as\), with

\[U_as-s\overline U_a=\Theta k_a.\]

This follows by applying \(1=K+\Theta\Lambda\) to \(U_as\). The
finite/global comparison is now fully determined, with no new choice of
arithmetic unit:

\[b_Z=\Lambda R:E\to V,\qquad R^s=s\sigma=R-\Theta b_Z.\tag{31}\]

Define \(c_a^s=k_a\sigma\). Substitution in (15) yields

\[
\boxed{c_a^s=c_a-U_ab_Z+b_Za^A.}\tag{32}
\]

To verify the equality in \(V\), use \(\overline U_a\sigma=\sigma a^A\)
from (18) to compute

\[
\begin{aligned}
U_aR^s-R^sa^A
&=U_aR-Ra^A-\Theta U_ab_Z+\Theta b_Za^A\\
&=\Theta(c_a-U_ab_Z+b_Za^A)=\Theta k_a\sigma.
\end{aligned}
\]

Theta injectivity proves (32). Its derivative is the exact infinitesimal
bridge

\[\boxed{k_D\sigma=\phi_*\ell-Db_Z+b_ZA,\qquad k_D=\Lambda Ds.}\tag{33}\]

All terms are defined continuous maps; differentiation is allowed by
Section 1. This proves that the explicitly evaluated Duhamel integral
and the transcript's global continuous cocycle represent the same actual
arithmetic extension, through the displayed coboundary.

The full extension class is unchanged. Applying \(h(D)\) to (31) gives

\[h(D)R^s(u)=\Theta\left(R_Z(u)(D)\phi_*-h(D)b_Zu\right).\]

Modulo \(h(D)V\), its source class equals that of \(R_Z(u)(D)\phi_*\).
Under \(j_hH\) its value is exactly \(\varepsilon u\), by (7). There is
no quotient operation that discards the nonconstant coefficients of
\(\varepsilon\).

At cochain level the two finite representative maps \(r^s\) and \(r\)
have the difference

\[r^s-r=dL^\pm,\qquad
L^+(u)=(-b_Zu,0),\qquad L^-(u)=(0,\widehat{b_Zu}).\tag{34}\]

The minus-face sign follows by applying \(d(0,\widehat b)=-\Theta b\).
Equations (25), (32), and (34) give the exact homotopy comparison

\[H_a^{s,\pm}=H_a^\pm+\mathcal U_a^0L^\pm-L^\pm a^A.\]

These identities retain the joint degree-zero difference as in (28),
rather than asserting strict agreement of the two face homotopies.

\subsection{7. Tensor powers with every cohomological sign and jet
coefficient}\label{proof04:tensor-powers-with-every-cohomological-sign-and-jet-coefficient}

Let \(k\geq1\). On the original ordered tensor complex \(C^{\otimes k}\)
use

\[
d(x_1\otimes\cdots\otimes x_k)
=\sum_{j=1}^k(-1)^{\sum_{i<j}|x_i|}
x_1\otimes\cdots\otimes dx_j\otimes\cdots\otimes x_k.
\tag{35}
\]

This is exactly the tau-base product resolution differential. It squares
to zero because the two orders of applying differentials in different
factors have opposite signs; a repeated differential in one factor is
zero.

Write \(f_a=U_aR:E\to C^1\), \(g_a=Ra^A:E\to C^1\), and
\(h_a=H_a^\pm:E\to C^0\) for any admitted choice of face in each factor.
The explicit degree-\(-1\) homotopy on the top packet is

\[
\begin{aligned}
\mathcal H_{a,k}(u_1\otimes\cdots\otimes u_k)
=\sum_{j=1}^k(-1)^{j-1}
f_au_1\otimes\cdots\otimes f_au_{j-1}
\otimes h_au_j\otimes
g_au_{j+1}\otimes\cdots\otimes g_au_k.
\end{aligned}\tag{36}
\]

In this evaluated formula the factor \((-1)^{j-1}\) is displayed exactly
once. If the same formula is instead written as a tensor of graded
linear maps, it is the Koszul evaluation sign of the degree-\(-1\) map
in slot \(j\) crossing the preceding \(j-1\) inputs, each of degree
\(1\); it must not be inserted a second time.

Only the slot containing \(h_au_j\) contributes to the differential,
because the other slots are already in top degree. Its differential sign
from (35) is \((-1)^{j-1}\), which cancels the prefactor. Since
\(dh_a=f_a-g_a\), the result is the telescoping identity

\[
\boxed{d\mathcal H_{a,k}
=(U_aR)^{\otimes k}-(Ra^A)^{\otimes k}.}\tag{37}
\]

For \(k=2\) this reads
\(\mathcal H_{a,2}=h_a\otimes g_a-f_a\otimes h_a\) on evaluated inputs,
and its differential is \((f_a-g_a)\otimes g_a+f_a\otimes(f_a-g_a)\).
This displays the original minus sign without relying on an implicit
convention.

Define

\[R_k=R^{\otimes k},\qquad
A_k=\sum_{j=1}^k1^{\otimes(j-1)}\otimes A\otimes1^{\otimes(k-j)},
\qquad T_{t,k}=T_t^{\otimes k}=e^{tA_k}.\]

For the ordered local block \((\rho_1,\ldots,\rho_k)\) the action is
exactly

\[
T_{t,k}
=e^{t\sum_i\rho_i}
\sum_{0\leq j_i<m_{\rho_i}}
\frac{t^{j_1+\cdots+j_k}}{j_1!\cdots j_k!}
N_{\rho_1}^{j_1}\otimes\cdots\otimes N_{\rho_k}^{j_k}.
\tag{38}
\]

The formula follows by multiplying (19); the nilpotents in different
factors commute, and no term has been discarded. The representative is
still the tensor of the maps using \(\varepsilon\) in (10), hence the
coefficient unit is \(\varepsilon^{\otimes k}\) in the ordered local
algebra.

For the multivariable Mellin transform \(\mathcal M_k\) with measure
\(\prod_i dx_i/x_i\), the top boundary \(B_{a,k}=d\mathcal H_{a,k}\)
has, on a pure coefficient tensor, the exact expression

\[
\begin{aligned}
\mathcal M_kB_{a,k}(\mathbf s)
={}&\prod_{i=1}^k\left[a^{s_i}v(s_i)R_Z(u_i)(s_i)\right]\\
&-\prod_{i=1}^k\left[v(s_i)R_Z(a^Au_i)(s_i)\right].
\end{aligned}\tag{39}
\]

All integrals factor by absolute convergence of the Schwartz Mellin
integrals. Applying the full packet jet in each factor gives zero, as
also follows from (37). The homotopy (36) exhibits the actual tensor
relation giving that zero. It is not inferred from a dimension count or
an assumed vanishing tensor kernel.

Every map in (36) is continuous. More explicitly, for seminorms
\(p_1,\ldots,p_k\) on the appropriate target factors, the projective
tensor seminorm of one summand on a pure input is bounded by the product
of the corresponding operator bounds from (1), (10), and (14), times
\(\prod_i\|u_i\|\). Sum over the \(k\) terms and use the defining
infimum for the projective seminorm on arbitrary finite sums. Thus all
maps extend to completed projective tensor products, and (37) persists
by density and continuity. The coefficient domain \(E^{\otimes k}\) is
finite dimensional, but this continuity statement identifies its targets
in the same completed topology used by the original global construction.

The support rule is applied to each admitted factor before the tensor
balancing quotient: an external absent factor is absorbing, while a zero
amplitude in an active factor remains its supported tensor zero. The
chosen face in that factor determines which of (25) occurs. Changing
that face replaces its homotopy by (28); substitution into (36) gives an
explicit tensor of the retained \(H^0\) cycle with the other degree-one
factors. No identification of external tensor support with a Cartesian
product of labels is needed or asserted.

This tensor homotopy difference can be detected without discarding a
possible lower-degree boundary. Define cochain maps \(\pi_0:C\to V[0]\)
by \(\pi_0^0(v,w)=w\), \(\pi_0^1=0\), and \(\pi_1:C\to Q[-1]\) by
\(\pi_1^0=0\), \(\pi_1^1=q\). In the difference between the all-plus and
all-minus versions of (36), apply \(\pi_0\) in slot \(j\) and \(\pi_1\)
in the others. Every summand but the \(j\)-th vanishes. The result is

\[(-1)^{j-1}(\sigma a^A)^{\otimes(j-1)}\otimes\widehat c_a
\otimes(\sigma a^A)^{\otimes(k-j)}.\]

It is a nonzero map whenever \(a\ne1\): \(\sigma a^A\) is injective by
(11), \(\widehat c_a\) is nonzero by Section 5, and a tensor of nonzero
vectors over \(\mathbb C\) is nonzero, as witnessed by a tensor of
linear functionals nonzero on those vectors. The target complex has zero
differential, so the detected class cannot be a boundary. This proves
the exact survival of the homotopy difference in degree \(k-1\), where
one factor carries the original \(H^0(C)\) and the others carry \(Q\).

\subsection{\texorpdfstring{8. The complete \(dx\) boundary pairing, its
derivative, and tensor
mass}{8. The complete dx boundary pairing, its derivative, and tensor mass}}\label{proof04:the-complete-dx-boundary-pairing-its-derivative-and-tensor-mass}

The inclusion \(\mathscr B\hookrightarrow L^2(\mathbb R_{>0},dx)\) is
continuous. For example on \((0,1)\) use \(|F(x)|\leq p_{-1,0}(F)x\),
and on \((1,\infty)\) use \(|F(x)|\leq p_{1,0}(F)x^{-1}\). Both squared
majorants are integrable. Use the inner product conjugate-linear in its
first argument. Direct substitution, retaining \(dx\), gives

\[\langle U_aF,U_aG\rangle=a\langle F,G\rangle.\tag{40}\]

Let \(G=R^*R\) in the retained coefficient basis and let
\(B_a=\Theta c_a:E\to L^2(dx)\). Since \(J_ZR=1\), \(G\) is positive
definite. Equation (15), \(U_aR=Ra^A+B_a\), gives the exact finite
boundary identity

\[
\boxed{
(a^A)^*Ga^A-aG
=-(U_aR)^*B_a-B_a^*(U_aR)+B_a^*B_a.
}\tag{41}
\]

On vectors \(u,w\), its right side is

\[-\langle U_aRu,\Theta c_aw\rangle
-\langle\Theta c_au,U_aRw\rangle
+\langle\Theta c_au,\Theta c_aw\rangle.\]

To check the signs, solve \(Ra^A=U_aR-B_a\) and expand its squared
pairing. Equation (40) supplies \(aG\). Equivalently the right side is
\(-(Ra^A)^*B_a-B_a^*Ra^A-B_a^*B_a\); these two forms agree by (15). Thus
the positive sign in the version using \(U_aR\) is forced by the
expansion, not a sign choice.

Every term is finite and depends smoothly on \(a\), by the Frechet
bounds already proved and the continuous Hilbert inclusion. For a full
generalized packet vector \(u\), the left side is evaluated using (19),
retaining every power of \(\log a\). For a true eigenvector
\(Au=\rho u\) it is \((a^{2\operatorname{Re}\rho}-a)\,u^*Gu\).

The infinitesimal identity obtained by differentiating (41) at
\(t=\log a=0\) is

\[
W=A^*G+GA-G
=-R^*f_0\ell-\ell^*f_0^*R.\tag{42}
\]

This agrees exactly with (12) and the original integration by parts.
Indeed for \(F,H\in\mathscr B\),

\[\langle DF,H\rangle+\langle F,DH\rangle
=-\int_0^\infty x(\overline FH)'dx
=\langle F,H\rangle,\]

since the endpoint term \([x\overline F(x)H(x)]_0^\infty\) vanishes by
the displayed weighted bounds. This independently verifies the mass term
\(-G\) and the two negative boundary cross terms in (42).

The finite defect is exactly the integrated infinitesimal one:

\[
\boxed{T_t^*GT_t-e^tG
=\int_0^t e^{t-r}T_r^*WT_r\,dr.}\tag{43}
\]

Set \(F(t)=T_t^*GT_t-e^tG\). Direct differentiation gives
\(F'(t)=F(t)+T_t^*WT_t\) and \(F(0)=0\). Multiplying by \(e^{-t}\) and
integrating proves (43) with the original orientation for \(t<0\). This
is a finite-dimensional exact matrix integration, valid with all the
nilpotents in (19).

For the global representative \(R^s\) from (31), replace \(R\) by
\(R^s\) and \(c_a\) by the fully computed \(c_a^s\) in (32). The finite
expansion (41) then holds with

\[
G^s=(R-\Theta b_Z)^*(R-\Theta b_Z),\quad
\Theta c_a^s=\Theta(c_a-U_ab_Z+b_Za^A).
\]

Its infinitesimal boundary is the changed map

\[
B_D^s=\Theta k_D\sigma
=\Theta(\phi_*\ell-Db_Z+b_ZA).
\tag{43a}
\]

Indeed differentiating the explicit \(c_{e^t}^s\) at \(t=0\) gives
\(k_D\sigma\) by (33), and direct differentiation of
\(R^s=R-\Theta b_Z\) gives \(DR^s-R^sA=B_D^s\) by (12). Apply the same
\(dx\) integration-by-parts identity to \(R^su,R^sw\) for arbitrary
coefficient vectors. Substituting \(DR^s=R^sA+B_D^s\) and moving the two
boundary terms to the other side proves exactly

\[
W^s=A^*G^s+G^sA-G^s
=-(R^s)^*B_D^s-(B_D^s)^*R^s.
\tag{43b}
\]

Thus the factor \(f_0\ell\) in (42) is replaced by \(B_D^s\), not
retained as the old infinitesimal source boundary. Finally
differentiating \(T_t^*G^sT_t-e^tG^s\) as in the proof of (43) gives
\(F_s'(t)=F_s(t)+T_t^*W^sT_t\), \(F_s(0)=0\), and therefore

\[T_t^*G^sT_t-e^tG^s=\int_0^t e^{t-r}T_r^*W^sT_r\,dr.\tag{43c}\]

This is exactly the restriction of the A1694 global pairing identity to
the actual packet \(\sigma E\). It proves the finite/global relationship
between both the cocycle and its boundary pairings, without assuming an
unproved vanishing of either cross term.

For \(k\) factors use the Hilbert pairing with measure
\(dx_1\cdots dx_k\). Fubini gives its value on pure tensors as the
product of the original pairings, hence

\[\langle U_a^{\otimes k}F,U_a^{\otimes k}H\rangle
=a^k\langle F,H\rangle.\]

Let \(G_k=G^{\otimes k}\), \(B_{a,k}\) be the actual boundary (37), and
\(\mathcal R_a=(U_aR)^{\otimes k}\). Then the exact tensor version of
(41) is

\[
\boxed{(a^A)^{*\otimes k}G_k(a^A)^{\otimes k}-a^kG_k
=-\mathcal R_a^*B_{a,k}-B_{a,k}^*\mathcal R_a+B_{a,k}^*B_{a,k}.}\tag{44}
\]

Indeed \(R_k(a^A)^{\otimes k}=\mathcal R_a-B_{a,k}\) by (37), and
expansion proves the formula. Its derivative has the exact ordered sum

\[
W_k=A_k^*G_k+G_kA_k-kG_k
=\sum_{j=1}^kG^{\otimes(j-1)}\otimes W\otimes G^{\otimes(k-j)},
\tag{45}
\]

because each of the \(k\) copies of the term \(-G\) in (42) contributes
one copy of \(-G_k\). The finite integral is

\[T_{t,k}^*G_kT_{t,k}-e^{kt}G_k
=\int_0^t e^{k(t-r)}T_{r,k}^*W_kT_{r,k}\,dr.\tag{46}\]

The proof differentiates the left side exactly as in (43), with mass
exponent \(k\). Equations (38) and (46) retain all tensor nilpotents.
They are actual formulas for the original representative boundary
contribution on the original tau-base tensor complex.

\clearpage

\subsection{9. Source pins and status of the
calculation}\label{proof04:source-pins-and-status-of-the-calculation}

The complete input files read for this calculation were:

\begin{itemize}
\tightlist
\item
  \texttt{arithmetic\_\allowbreak{}input.tex} (original September 12 cumulative
  source), SHA256
  \nolinkurl{a50b0fb587b644c4ea94dba45b2b423d038f2bfbe7fef188543939ba850e58d2}.
\item
  \emph{Tau\_Base\_Cohomology}, \texttt{NOTE.md}, SHA256
  \nolinkurl{d03e71ad18f187bd89880cb8ca6fc9c5d6f260b3de8e8e5702c27db97e7b63c3}.
\item
  Original transcript turn \texttt{A1694.md}, SHA256
  \nolinkurl{a87bff4960c0a1742734e9e711780e0533ea76ba6f46ceab1fea115a2fd66f18};
  its original locator is recorded in Section 0.
\item
  \emph{Tau\_Global\_Comparison}, \texttt{NOTE.tex}, SHA256
  \nolinkurl{50284fcebae0dc859ac398a00839a29e15be60c008e02314ac1f35c31950ed01};
  its equations are cited at the beginning.
\end{itemize}

The completed new formulas are (13)--(24), their exact finite/global
comparison (31)--(34), the evaluated tensor homotopy and jets
(36)--(39), and the integrated finite and tensor boundary expressions
(43)--(46). Sections 1, 2, 5, 6, and the unintegrated boundary pairing
explicitly reprove the source interfaces needed to type and verify those
formulas. The global retraction, global cocycle, and support comparison
in A1694 remain credited to that earlier calculation.

The resulting boundary terms have been computed, including their
orientations, topologies, full arithmetic coefficients, and product
masses. None has been set to zero by calling it a theta boundary. No
estimate forcing the sign or vanishing of the original global weight
defect is asserted by these identities. That precise limit on the
conclusion does not delete the displayed maps, the full quotient \(Q\),
or the original tau-base construction.
